% Generated draft: requires mathematical and source review.
\documentclass[11pt]{article}
\usepackage[margin=25mm]{geometry}
\usepackage{fontspec}
\setmainfont{Latin Modern Roman}
\usepackage{amsmath,amssymb,mathtools,unicode-math}
\setmathfont{Latin Modern Math}
\usepackage{enumitem,needspace,xurl,hyperref}
\setlength{\parindent}{0pt}
\setlength{\parskip}{5pt}
\setlength{\emergencystretch}{3em}
\newcommand{\N}{\mathbb{N}}
\newcommand{\Z}{\mathbb{Z}}
\newcommand{\Q}{\mathbb{Q}}
\newcommand{\R}{\mathbb{R}}
\newcommand{\C}{\mathbb{C}}
\newcommand{\F}{\mathbb{F}}
\newcommand{\A}{\mathbb{A}}
\newcommand{\PP}{\mathbb{P}}
\newcommand{\E}{\mathbb{E}}
\newcommand{\eps}{\varepsilon}
\newcommand{\dd}{\,\mathrm{d}}
\begin{document}
\section*{UnsolvedMath ID: 1101302}
\textbf{Source problem code:} catalogue entry
\subsection{Definitions and mathematical statement}
Let $n\geq0$ be an integer. Use integral Poincar\'e duality: a PD$(n)$ group $G$ has cohomological dimension $n$, the trivial $\Z G$-module $\Z$ admits a finite resolution by finitely generated projective modules, and
\[
H^i(G;\Z G)=0\quad(i\ne n),\qquad
H^n(G;\Z G)\cong\Z\quad\text{as an abelian group}.
\]
The action on the top cohomology may be the orientation-character action; orientability is not imposed.

A group is solvable if its derived series $G^{(0)}=G$, $G^{(j+1)}=[G^{(j)},G^{(j)}]$ reaches the trivial group after finitely many steps. It is polycyclic if it has a finite subnormal series with cyclic factors. Its Hirsch length, when polycyclic, is the number of infinite cyclic factors in such a series.

Kevin Whyte's question asks whether
\[
\forall n\geq0\ \forall G,\qquad
\bigl(G\text{ solvable and PD}(n)\bigr)
\Longrightarrow G\text{ polycyclic}.
\]
The supplied research summary reports an affirmative theorem of Bieri and the sharper classification: the solvable integral PD$(n)$ groups are precisely the torsion-free polycyclic groups of Hirsch length $n$. This is a reported literature resolution, not a proof supplied here.
\subsection{Short English statement}
The question asks whether integral Poincar\'e duality forces a solvable group to be polycyclic. The supplied research summary attributes an affirmative answer to Bieri, with dimension equal to the Hirsch length.
\subsection{Sources}
\begin{itemize}
\item Local UnsolvedMath registry, permanent ID 1101302, exact source code catalogue entry, ``Questions in Geometric Group Theory --- Q 13.2''; status solved, classification SOLVED-IN-LITERATURE.
\item Bestvina, \emph{Questions in Geometric Group Theory} (2004), Question 13.2, PDF page 24, attributed to Kevin Whyte, \url{https://www.math.utah.edu/~bestvina/eprints/questions-updated.pdf}. The record identifies manual transcription of a public author-hosted PDF, accessed July 29, 2026; the source presents it as open and current status was not independently updated at extraction.
\item Bieri (1972), and Theorem 9.23 of Bieri's book (1976/1981), as cited by the supplied research summary. The summary attributes the affirmative theorem and torsion-free polycyclic classification to these sources and says the Bestvina entry was apparently outdated already in 2004. Titles and fuller bibliographic details are not supplied.
\end{itemize}
\end{document}
